\subsection{无界算子的若干谱分析结果}
\label{sec:chap2-unbounded-operator-spectral-analysis}
\renewcommand{\thestatement}{6.\number\numexpr\value{statement}-93\relax}

为把有限维谱理论推广到无限维, 考虑从 Hilbert 空间 $H$ 到自身、但只在子空间 $D(A)$ 上定义的算子 $A$. 后文将其应用于 Stokes 算子.

\subsubsection{定义}
\label{subsec:chap2-unbounded-operator-definitions}

无界算子由 Hilbert 空间 $H$、线性子空间 $D(A)\subset H$ 以及线性映射 $A:D(A)\to H$ 给出. 本书假设 $D(A)$ 在 $H$ 中稠密, 且图像
\[
 G(A)=\{(u,Au):u\in D(A)\}
\]
在 $H\times H$ 中闭.

\par\noindent\textbf{Remark 6.1:}\quad 若闭算子 $A$ 的稠密定义域 $D(A)\neq H$, 则不存在 $C>0$ 使 $\|Au\|_H\leq C\|u\|_H$ 对所有 $u\in D(A)$ 成立. 否则由稠密性和图像闭性可推出每个 $u\in H$ 都属于 $D(A)$.

\begin{Definition}
\label{def:chap2-adjoint-unbounded-operator}
设 $A:D(A)\subset H\to H$ 有稠密定义域. 定义
\[
 D(A^*)=\{u\in H:v\in D(A)\mapsto(Av,u)_H
 \text{ 关于 }H\text{ 范数连续}\}.
\]
对 $u\in D(A^*)$, 唯一元素 $A^*u\in H$ 由
\[
 (Av,u)_H=(v,A^*u)_H,
 \qquad\forall v\in D(A)
\]
确定. 称 $A^*$ 为 $A$ 的伴随算子.
\end{Definition}

\begin{Definition}
\label{def:chap2-self-adjoint-unbounded-operator}
若 $D(A^*)=D(A)$ 且 $A^*u=Au$ 对所有 $u\in D(A)$ 成立, 则称 $A$ 自伴.
\end{Definition}

\begin{Proposition}
\label{prop:chap2-inverse-self-adjoint}
设无界算子 $A$ 自伴, $A:D(A)\to H$ 为双射, 且 $A^{-1}:H\to H$ 连续. 则有界算子 $A^{-1}$ 自伴.
\end{Proposition}

\begin{Proof}
连续性给出 $D((A^{-1})^*)=H$. 对 $u,v\in H$, $A^{-1}u,A^{-1}v\in D(A)$, 且
\[
 (u,A^{-1}v)_H=(A(A^{-1}u),A^{-1}v)_H
 =(A^{-1}u,A(A^{-1}v))_H=(A^{-1}u,v)_H.
\]
\end{Proof}

\begin{Proposition}
\label{prop:chap2-graph-norm-hilbert}
设 $A$ 为 $H$ 中定义域为 $D(A)$ 的闭无界算子. 赋予 $D(A)$ 内积
\begin{equation}
 (u,v)_{D(A)}=(u,v)_H+(Au,Av)_H.
 \label{eq:chap2-graph-inner-product}
\end{equation}
则 $D(A)$ 是 Hilbert 空间, $D(A)\hookrightarrow H$ 连续, 且 $A:D(A)\to H$ 连续.
\end{Proposition}

\par\noindent\textbf{Remark 6.2:}\quad 上述范数等价于图范数 $\|u\|_H+\|Au\|_H$.

\subsubsection{谱理论的初等结果}
\label{subsec:chap2-elementary-spectral-theory}

\begin{Theorem}[Compact self-adjoint operators]
\label{thm:chap2-compact-self-adjoint-spectral}
设 $H$ 为可分无限维 Hilbert 空间, $T:H\to H$ 为紧自伴有界算子. 则 $H$ 有一组由 $T$ 的特征向量组成的标准正交基. 其特征值均为实数, 可排列成趋于 $0$ 的序列.
\end{Theorem}

\begin{Theorem}[Operators with compact inverse]
\label{thm:chap2-compact-inverse-spectral}
设 $H$ 为可分无限维 Hilbert 空间, $A:D(A)\subset H\to H$ 自伴, 从 $D(A)$ 到 $H$ 为双射, 且 $D(A)\hookrightarrow H$ 紧. 则存在 $H$ 的标准正交基 $(w_k)_{k\geq1}$ 以及实数 $\lambda_k$, 使
\[
 w_k\in D(A),\qquad Aw_k=\lambda_kw_k,
\]
并可使 $(|\lambda_k|)_k$ 单调递增且趋于 $+\infty$. 此外, $(w_k)_k$ 在 $D(A)$ 中为完备正交族.
\end{Theorem}

\begin{Proof}
由 Proposition~\ref{prop:chap2-inverse-self-adjoint}, $A^{-1}$ 是紧自伴有界算子. Theorem~\ref{thm:chap2-compact-self-adjoint-spectral} 给出 $H$ 的标准正交基 $(w_k)_k$ 以及非零特征值 $\mu_k\to0$. 令 $\lambda_k=1/\mu_k$, 则 $Aw_k=\lambda_kw_k$ 且 $|\lambda_k|\to+\infty$.

对 $k\neq l$,
\[
 (w_k,w_l)_{D(A)}=(1+\lambda_k\lambda_l)(w_k,w_l)_H=0.
\]
若 $u\in D(A)$ 与所有 $w_k$ 在 $D(A)$ 中正交, 则
\[
 0=(u,w_k)_{D(A)}=(1+\lambda_k^2)(u,w_k)_H,
\]
故 $u=0$, 从而该族在 $D(A)$ 中完备.
\end{Proof}

\begin{Proposition}
\label{prop:chap2-domain-spectral-characterization}
在 Theorem~\ref{thm:chap2-compact-inverse-spectral} 的假设下,
\[
 D(A)=\left\{u\in H:\sum_{k\geq1}\lambda_k^2(u,w_k)_H^2<+\infty\right\}.
\]
\end{Proposition}

\begin{Proof}
$(w_k/\|w_k\|_{D(A)})_k$ 是 $D(A)$ 的标准正交基, 且
\[
 (u,w_k)_{D(A)}=(1+\lambda_k^2)(u,w_k)_H,
 \qquad\|w_k\|_{D(A)}^2=1+\lambda_k^2.
\]
Parseval 恒等式给出必要性. 反之, 若加权级数收敛, 则 $\sum_k(u,w_k)_Hw_k$ 在 $D(A)$ 中收敛到某个 $\widetilde u$. 连续嵌入说明它在 $H$ 中也收敛, 标准正交基的完备性给出 $u=\widetilde u\in D(A)$.
\end{Proof}

以下假设 $A$ 非负. 对 $s\geq0$, 定义
\[
 D(A^s)=\left\{u\in H:\sum_{k\geq1}\lambda_k^{2s}(u,w_k)_H^2<+\infty\right\},
 \qquad A^su=\sum_{k\geq1}\lambda_k^s(u,w_k)_Hw_k.
\]

\begin{Proposition}
\label{prop:chap2-fractional-power-properties}
有:
\begin{enumerate}
\item $A^1=A$, $D(A^0)=H$, 且 $A^0$ 为恒等算子.
\item 对 $s>0$, $D(A^s)$ 为 Hilbert 空间, $A^s$ 是从 $D(A^s)$ 到 $H$ 的非负自伴同构, $(w_k)_k$ 在 $D(A^s)$ 中完备.
\item 若 $0\leq s<s'$, 则 $D(A^{s'})\subsetneq D(A^s)$, 且该嵌入紧.
\end{enumerate}
\end{Proposition}

\begin{Proof}
只证紧性. 若 $(u_j)_j$ 在 $D(A^{s'})$ 中有界, 则对充分大的 $n_0$,
\begin{equation}
 \sum_{n\geq n_0}\lambda_n^{2s}|(u_j,w_n)_H|^2\leq\varepsilon^2/2,
 \label{eq:chap2-fractional-power-tail}
\end{equation}
对所有 $j$ 一致成立. 前 $n_0$ 个坐标位于有限维有界集, 可由有限个小球覆盖, 即
\begin{equation}
 \sum_{n<n_0}\lambda_n^{2s}|(u_j-v^i,w_n)_H|^2\leq\varepsilon^2/2
 \label{eq:chap2-fractional-power-head}
\end{equation}
对某个有限集合中的 $v^i$ 成立. 两式相加得到 $D(A^s)$ 中的有限 $\varepsilon$-覆盖.
\end{Proof}

对 $s<0$, 以范数 $\|u\|_{D(A^s)}^2=\sum_k\lambda_k^{2s}(u,w_k)_H^2$ 完备化 $H$, 得到大于 $H$ 的 Hilbert 空间 $D(A^s)$.

\begin{Proposition}
\label{prop:chap2-negative-power-duality}
对 $s\geq0$, 把 $H$ 与其对偶等同后,
\[
 D(A^s)'\simeq D(A^{-s}),
\]
且
\[
 \langle u,v\rangle_{D(A^{-s}),D(A^s)}
 =\sum_{k\geq1}(u,w_k)_H(v,w_k)_H.
\]
\end{Proposition}

\subsubsection{半群理论的应用}
\label{subsec:chap2-semigroup-applications}

设 $A$ 满足上述自伴、双射、非负和紧逆条件. 考虑
\begin{equation}
 \begin{cases}
 \dfrac{\mathrm{d}u}{\mathrm{d}t}+Au=0,\\
 u(0)=u_0.
 \end{cases}
 \label{eq:chap2-homogeneous-semigroup-problem}
\end{equation}

\begin{Proposition}[Definition and proposition]
\label{prop:chap2-semigroup-definition}
若 $u_0=\sum_{k\geq1}u_{0,k}w_k\in H$, 则对 $t\geq0$ 可定义
\[
 e^{-tA}u_0=\sum_{k\geq1}u_{0,k}e^{-t\lambda_k}w_k\in H.
\]
对任意 $s\geq0$ 和 $t>0$, $e^{-tA}u_0\in D(A^s)$.
\end{Proposition}

\par\noindent\textbf{Remark 6.3:}\quad $e^{-(t+s)A}=e^{-tA}e^{-sA}=e^{-sA}e^{-tA}$, 因而 $(e^{-tA})_{t\geq0}$ 称为与 $-A$ 关联的半群.

\begin{Proof}
因 $\lambda_k,t\geq0$, 级数在 $H$ 中有定义. 对 $s>0$,
\[
 \lambda_k^se^{-t\lambda_k}=(t\lambda_k)^se^{-t\lambda_k}t^{-s}
 \leq C_st^{-s}.
\]
故
\[
 \sum_{k\geq1}(\lambda_k^su_{0,k}e^{-t\lambda_k})^2
 \leq C_s^2t^{-2s}\sum_{k\geq1}|u_{0,k}|^2<+\infty.
\]
\end{Proof}

\begin{Theorem}
\label{thm:chap2-homogeneous-semigroup-solution}
对每个 $u_0\in H$, 问题\eqref{eq:chap2-homogeneous-semigroup-problem}有唯一解
\[
 u\in C^0([0,+\infty),H)\cap C^1((0,+\infty),D(A)),
\]
且
\begin{equation}
 u(t)=e^{-tA}u_0.
 \label{eq:chap2-homogeneous-semigroup-formula}
\end{equation}
\end{Theorem}

\begin{Proof}
若 $u_0=0$, 对 $0<\varepsilon<T$ 应用 Theorem~\ref{thm:chap2-lions-magenes-product},
\[
 \|u(T)\|_H^2-\|u(\varepsilon)\|_H^2
 =-2\int_\varepsilon^T(Au(t),u(t))_H\,\mathrm{d}t\leq0.
\]
令 $\varepsilon\to0$ 得 $u(T)=0$, 故解唯一. 式\eqref{eq:chap2-homogeneous-semigroup-formula}定义的级数具有所述正则性, 且对 $t>0$ 可逐项求导:
\[
 \frac{\mathrm{d}u}{\mathrm{d}t}
 =-\sum_{k\geq1}\lambda_ku_{0,k}e^{-t\lambda_k}w_k=-Au(t).
\]
\end{Proof}

\begin{Theorem}
\label{thm:chap2-duhamel-semigroup}
设 $u_0\in H$, $f\in C^1([0,+\infty),H)$. 则问题
\begin{equation}
 \begin{cases}
 \dfrac{\mathrm{d}u}{\mathrm{d}t}+Au=f,\\
 u(0)=u_0
 \end{cases}
 \label{eq:chap2-inhomogeneous-semigroup-problem}
\end{equation}
有唯一解 $u\in C^0([0,+\infty),H)\cap C^1((0,+\infty),H)\cap C^0((0,+\infty),D(A))$, 且
\begin{equation}
 u(t)=e^{-tA}u_0+\int_0^te^{-(t-s)A}f(s)\,\mathrm{d}s.
 \label{eq:chap2-duhamel-formula}
\end{equation}
\end{Theorem}

\begin{Proof}
在式\eqref{eq:chap2-duhamel-formula}的积分中作变量代换 $s\mapsto t-s$, 得
\[
 u(t)=e^{-tA}u_0+\int_0^te^{-sA}f(t-s)\,\mathrm{d}s.
\]
$f$ 为 $C^1$ 的假设允许逐项求导, 再分部积分即得方程和初值. 唯一性由齐次问题得到.
\end{Proof}

\par\noindent\textbf{Remark 6.4:}\quad 式\eqref{eq:chap2-duhamel-formula}对正则性较低的源项仍成立, 但需相应减弱解的概念.
